% Generated draft: requires mathematical and source review.
\documentclass[11pt]{article}
\usepackage[margin=25mm]{geometry}
\usepackage{fontspec}
\setmainfont{Latin Modern Roman}
\usepackage{amsmath,amssymb,mathtools,unicode-math}
\setmathfont{Latin Modern Math}
\usepackage{enumitem,needspace,xurl,hyperref}
\setlength{\parindent}{0pt}
\setlength{\parskip}{5pt}
\setlength{\emergencystretch}{3em}
\newcommand{\N}{\mathbb{N}}
\newcommand{\Z}{\mathbb{Z}}
\newcommand{\Q}{\mathbb{Q}}
\newcommand{\R}{\mathbb{R}}
\newcommand{\C}{\mathbb{C}}
\newcommand{\F}{\mathbb{F}}
\newcommand{\A}{\mathbb{A}}
\newcommand{\PP}{\mathbb{P}}
\newcommand{\E}{\mathbb{E}}
\newcommand{\eps}{\varepsilon}
\newcommand{\dd}{\,\mathrm{d}}
\begin{document}
\section*{UnsolvedMath ID: 30006627}
\textbf{Source problem code:} catalogue entry
\subsection{Definitions and mathematical statement}
Equip \(\R^3\) with its standard unit cubulation. Its lattice graph has vertex set \(\Z^3\), with an edge between two vertices precisely when their difference is one of the signed coordinate unit vectors. Let \(G\) be a finite connected subgraph of this graph; \(G\) need not be an induced subgraph. Define \(X(G)\) to contain the vertices and edges of \(G\), every standard unit square whose four vertices and four edges belong to \(G\), and every standard unit three-dimensional cube whose eight vertices and twelve edges belong to \(G\). Thus the one-skeleton of \(X(G)\) is \(G\).

For \(S\subseteq V(G)\), define the induced cubical subcomplex by
\[
X(G)[S]=\{\sigma\in X(G):\text{ every vertex of }\sigma\text{ belongs to }S\}.
\]
A cell of a cubical complex is maximal if it is not properly contained in another cell of that complex. A corner is a vertex contained in exactly one maximal cell; maximal cells here may have dimension zero, one, two, or three. A corner peeling is an ordering \(v_1,\ldots,v_n\) of all vertices of \(G\), where \(n=|V(G)|\), such that, for every \(1\leq i\leq n\), the vertex \(v_i\) is a corner of
\[
X(G)[\{v_1,\ldots,v_i\}].
\]
Equivalently, the vertices can be deleted as corners in the order \(v_n,\ldots,v_1\).

The supplied question assumes that \(X(G)\) is simply connected and that every connected component of its intersection with a coordinate hyperplane is simply connected. There is an unresolved specification in this hypothesis: the supplied record does not define the intended family of coordinate hyperplanes. To display this dependence explicitly, put
\[
H_{j,t}=\{x\in\R^3:x_j=t\},\qquad j\in\{1,2,3\},\quad t\in\R.
\]
For a specified family \(\mathcal H\) of such planes, the question has the following precise form:
\begin{quote}
For every finite connected lattice subgraph \(G\) as above, if \(X(G)\) is simply connected and every connected component of \(X(G)\cap H\) is simply connected for every \(H\in\mathcal H\), must \(G\) admit a corner peeling?
\end{quote}
Intersections are taken in the geometric realization in \(\R^3\). Empty intersections have no connected components and impose no condition. The record does not resolve whether \(\mathcal H\) means just the three planes \(H_{j,0}\), all integer translates, all real translates, or another specified family. The displayed question is therefore a formulation conditional on that missing choice, not a resolution of it. The supplied original excerpt ends before the conjectural assertion, so it cannot settle this qualification.
\subsection{Short English statement}
The question asks whether a finite connected subgraph of the three-dimensional integer lattice can be reduced by repeatedly deleting a vertex lying in exactly one maximal cube, square, edge, or isolated vertex, assuming its associated cubical complex and all prescribed plane-section components are simply connected. The intended family of planes needs clarification. The record labels the question open; this draft does not independently verify that status.
\subsection{Sources}
\begin{itemize}
\item Supplied local UnsolvedMath registry record, permanent ID \(30006627\), title: \emph{Corner Peeling in Cubical Subcomplexes of Integer Three-Space}; exact source problem code: \texttt{catalogue entry}.
\item Supplied attribution: Problem 15, \emph{Corners in subcomplexes of \(\Z^3\)}, V. Chepoi, from discussions with J. Chalopin and M. Kokkou. The supplied original excerpt defines the complex and corners but terminates at the incomplete text ``A'' before the conjectural statement.
\item Supplied report citation: \emph{Median Geometry and Applications} (2026), Oberwolfach Reports, DOI 10.4171/owr/2026/8, pp. 487--540. \url{https://doi.org/10.4171/owr/2026/8}.
\item Supplied literature source, labelled ``OWR report'': \url{https://ems.press/content/serial-article-files/53603}.
\item Supplied literature source, labelled ``related 2026 preprint'': \url{https://arxiv.org/abs/2602.12894}. No title or authors for this preprint are supplied in the record.
\item Literature assessment dated exactly \texttt{2026-08-22}: the supplied curation labels the assertion open and states that the cited 2026 work proves a theorem under an additional condition, rather than proving the conjecture. This is supplied dated curation, not an independently verified literature assessment; the additional condition is not specified in the record.
\item The registry marks the statement \texttt{corrected\_verified} and reports that the official report was checked, the question isolated, and notation repaired. These are supplied source-verification claims, not verification performed for this draft. The background describes the report as documenting workshop discussions and explicitly requests independent verification of time-sensitive status and cited claims.
\end{itemize}
\end{document}
